\begin{proof}[Proof of Theorem~\ref{I:U:thm:fixed}]
Negative thresholds are immediately rejected.  We give the proof for
$k\ge1$; for $k=0$ the sole competitor is zero and the same calculation
below has no candidate variables.

\paragraph{A fixed number of candidate variables.}
By Theorem~\ref{I:U:thm:attainment} and Lemma~\ref{I:U:lem:stable}, the inequality
$\delta_k(H_f)\le t$ is equivalent to the existence of a stable
$k$-state realization $g(w)=c^\top T_wb$ with error at most $t$.
A smaller stable realization can be padded with zero states.
Introduce the entries of $T_1,\ldots,T_d,b,c$, and a symmetric
$k\times k$ matrix $X$ as variables.  Impose
\begin{equation}\label{I:U:eq:fixed-stability}
 X\succ0,\qquad X-\sum_{a\in\Sigma}T_aXT_a^\top=\UId_k.
\end{equation}
These conditions are equivalent to mean-square stability of the tuple
$T$.  For the forward implication, positivity and
$\UId_k\succeq \Unorm{X}^{-1}X$ give a strict order contraction in the
order-unit norm determined by $X$; the spectral radius is therefore less
than one.  For the reverse implication, take
$X=\sum_{j\ge0}\Phi_T^j(\UId_k)$.
Positive definiteness in~\eqref{I:U:eq:fixed-stability} is expressed by
positivity of the $k$ leading principal minors of $X$.
The number of real variables is
\begin{equation}\label{I:U:eq:fixed-varcount}
 p=dk^2+2k+\frac{k(k+1)}2,
\end{equation}
which is independent of $n$.

\paragraph{Eliminating the joint Gramians.}
Let $m=n+k$, and form
\[
 B_a=\begin{pmatrix}A_a&0\\0&T_a\end{pmatrix},\qquad
 \beta=\begin{pmatrix}b_f\\b\end{pmatrix},\qquad
 \gamma=\begin{pmatrix}c_f\\-c\end{pmatrix}.
\]
The joint tuple $B$ is stable: a positive definite Stein solution for
the source and the matrix $X$ in~\eqref{I:U:eq:fixed-stability}, placed on
the diagonal, give a positive definite Stein solution for $B$.
Hence its exact reachability and observability Gramians are uniquely
defined by
\[
 P=\beta\beta^\top+\sum_aB_aPB_a^\top,\qquad
 Q=\gamma\gamma^\top+\sum_aB_a^\top QB_a.
\]
These are the usual finite-dimensional Gramian equations for weighted
automata; see also~\cite{sva2019}.
Using column vectorization, put
\begin{align}
 \mathcal L&=\UId_{m^2}-\sum_a B_a\otimes B_a,
 &\Delta&=\det\mathcal L,\label{I:U:eq:fixed-L}\\
 \widehat P
 &=\operatorname{unvec}\!\left(
       \operatorname{adj}(\mathcal L)\operatorname{vec}(\beta\beta^\top)
     \right),&&\nonumber\\
 \widehat Q
 &=\operatorname{unvec}\!\left(
       \operatorname{adj}(\mathcal L^\top)\operatorname{vec}(\gamma\gamma^\top)
     \right).&&
 \label{I:U:eq:fixed-adj}
\end{align}
Stability guarantees $\Delta\ne0$ throughout the feasible set, and
\begin{equation}\label{I:U:eq:fixed-rational-grams}
 P=\widehat P/\Delta,\qquad Q=\widehat Q/\Delta.
\end{equation}
Thus the Gramians are rational functions of the fixed collection of
variables; they need not be introduced as additional variables.

\paragraph{A polynomial norm constraint without square roots.}
The squared norm of the error Hankel is
$\lambda_{\max}(P^{1/2}QP^{1/2})$.
The matrices $PQ$ and $P^{1/2}QP^{1/2}$ have the same characteristic
polynomial: apply
$\det(z\UId-UV)=\det(z\UId-VU)$ with
$U=P^{1/2}$ and $V=P^{1/2}Q$.
In particular, all eigenvalues of $PQ$ are real and nonnegative,
including when $P$ is singular.
Define the polynomial matrix
\begin{equation}\label{I:U:eq:fixed-R}
 R=t^2\Delta^2\UId_m-\widehat P\widehat Q
   =\Delta^2(t^2\UId_m-PQ).
\end{equation}
Its eigenvalues are real.  Since $\Delta^2>0$, the error norm is at
most $t$ if and only if these eigenvalues are nonnegative.  By
Lemma~\ref{I:U:lem:coefftest}, this is equivalent to
\begin{equation}\label{I:U:eq:fixed-coeff}
 \det(z\UId_m+R)=z^m+a_1z^{m-1}+\cdots+a_m,
 \qquad a_1\ge0,\ldots,a_m\ge0.
\end{equation}
Each $a_i$ is a polynomial in the variables counted
in~\eqref{I:U:eq:fixed-varcount}.  Squaring the common denominator avoids
any assumption about its sign.  No matrix square root, eigenvalue
variable, or list of exponentially many principal minors is used.

\paragraph{Degrees, coefficient sizes, and construction time.}
Each entry of the $m^2\times m^2$ matrix $\mathcal L$ has degree at most two.
Consequently
\[
 \deg\Delta\le2m^2,\qquad
 \deg\widehat P,\deg\widehat Q\le2m^2,\qquad
 \deg a_i\le4m^3,
\]
where $t$ is the fixed rational input threshold.
Together,~\eqref{I:U:eq:fixed-stability} and~\eqref{I:U:eq:fixed-coeff}
give $O(k^2+m)$ polynomial equalities and inequalities of degree
$O(m^3)$ in $p$ variables.

Their coefficients have bit lengths polynomial in $m$ and the binary
input length.  To see this, clear the input denominators and use the
determinant expansion: each monomial coefficient is a sum of products
of polynomially many input coefficients, and the logarithm of the
number of summands is polynomial in $m$.  The same bound applies to
adjugate entries, their products, and the final characteristic
polynomial.  This argument bounds bit length; it does not call for
enumerating the determinant expansion.

For completeness, the polynomials themselves can be constructed in
polynomial bit time when $p$ is fixed.  Their degree bounds imply only
$m^{O(p)}$ possible monomials.  Evaluate determinants, adjugates, and
characteristic polynomials over the rationals on an integer
interpolation grid of polynomial side length, and perform exact
multivariate interpolation.  The grid contains polynomially many
points for fixed $p$, and every evaluation has polynomial bit
complexity.  Adjugates are evaluated directly by cofactors or a
division-free matrix algorithm, so singular interpolation points cause
no difficulty.  Their intermediate and output bit lengths are
polynomial by the same determinant bounds.  The interpolation step
also has polynomial bit complexity.

\paragraph{Decision.}
Real algebraic feasibility with a fixed number of variables is
decidable in time polynomial in the number of polynomials, their
degrees, and their coefficient bit lengths
\cite{renegar1992,basu2006}.  Apply this result to the formula above.
The source dimension enters through polynomial degree and coefficient
size, while the number of variables is the fixed number $p$.
This proves the claimed bit-complexity bound.
\end{proof}
